\section{Temporal Properties} \label{sec:temporal}

Before modeling the spectra we characterize each pulse in the time domain,
using exactly the same approved detector sets, background models, and source
windows as the spectral analysis (\S\ref{sec:bkg}). This serves two purposes.
First, the temporal properties---duration, pulse shape, minimum variability
timescale, and spectral lag---are physical observables in their own right,
probing the emission-region size, the outflow Lorentz factor, and the energy
dependence of the emitting sites. Second, the fitted pulse shape assigns every
Bayesian block an unambiguous phase (rise, peak, or decay), which the
spectral census of \S\ref{sec:obs} uses to ask the central question of this
survey: \emph{what is the spectral shape during each phase of a single
pulse?}

\subsection{Durations and Pulse Shapes} \label{sec:durations}

Durations ($T_{90}$, $T_{50}$) are measured from the background-subtracted,
detector-summed light curves of the approved NaI detectors, with
uncertainties from bootstrap resampling of the counts. Because the sample is
selected for isolated single-pulse morphology, the cumulative-counts method
is well behaved and the durations are insensitive to the exact integration
bounds. \PHt{$T_{90}$ distribution: median and IQR pending the full-sample
temporal run.}

Each pulse is then fitted with three shape models: the
\citet{Norris2005} exponential rise--decay profile, the
\citet{Kocevski2003} flux-normalized profile, and the double-hyperbolic-tangent
profile of \citet{Gowri2025}. The three parameterizations bracket the
freedom commonly used in the literature; we retain the best-fitting model
per burst (reduced $\chi^{2}$) and record the rise time, decay time, and
asymmetry $\phi = t_{\rm rise}/t_{\rm decay}$ with covariance-propagated
uncertainties. \PHt{Asymmetry distribution and best-model census pending.}
The fitted profile, not the raw light curve, defines the rise/peak/decay
phase boundaries used to tag Bayesian blocks, making the phase assignment
reproducible and model-explicit.

\subsection{Minimum Variability Timescale} \label{sec:mvt}

The minimum variability timescale (MVT) is computed with the Haar
structure-function method of \citet{Golkhou2014,Golkhou2015}, as implemented
and validated in the \textsc{MVTfermi} framework of \citet{Bala2025}. That
work establishes, from an extensive simulation suite, that the reliability
of a Haar MVT measurement depends jointly on the MVT value and its
signal-to-noise ratio ($\mathrm{SNR}_{\rm MVT}$), and provides an empirical
\emph{MVT Validation Curve} that classifies each measurement as a robust
detection or an upper limit. We adopt that classification scheme unchanged:
every MVT we quote carries its validation label, and values classified as
limits are never used as detections in correlations. Uncertainties are
propagated with a Monte Carlo over Poisson realizations of the raw counts
with the fitted background model subtracted. \PHt{Full-band MVT
distribution, and energy-resolved MVT(E) in the standard sub-bands
(8--50, 50--300, 300--900 keV), pending the batch run.}

\subsection{Spectral Lags} \label{sec:lags}

Spectral lags between soft and hard energy bands are measured with the
cross-correlation method \citep{Norris1996}, with the peak of the CCF
located by a polynomial fit and its uncertainty estimated from Monte Carlo
resampling of the light curves. A lag is accepted only when the CCF peak is
significant against a bootstrap null distribution; below that threshold we
report an upper bound rather than a value. \PHt{Lag distribution and
lag--asymmetry / lag--MVT relations pending.}

Together these measurements give each burst a compact temporal fingerprint
(duration, shape, asymmetry, MVT, lag) obtained from the same gated
selections as the spectroscopy, so that every spectral statement in the
following sections can be conditioned on temporal phase without any
additional data processing.
